\chapter{Fundamentals}
\label{ch:fundamentals}

Description of the context of the work with all areas, methods, theories,
findings and technologies affected by the problem statement.

\section{Context}
\label{sec:context}

\subsection{Domain}

\subsection{Technologies}

An \gls{API} exposes functionality to other programs. Modern systems
increasingly combine an \acrfull{llm} with \acrshort{rag} to ground generated
output in a document corpus. Where the corpus is editorially maintained, a
\gls{CMS} typically owns it, while \gls{ML} components consume it downstream.

Note that only glossary entries actually referenced in the text appear in the
printed glossary --- unused entries in \texttt{glossary.tex} stay silent.

\subsection{Methods and Concepts}

\begin{definition}[Working definition]
\label{def:example}
A formal definition of a term introduced by this thesis. Theorem-like
environments are numbered per chapter.
\end{definition}
